\documentclass{article}
\usepackage{../../Self_Style}

\title{Phys 115B HW 5}
\author{Zih-Yu Hsieh}
\date{\today}

\begin{document}
\maketitle

\section*{1 D}
\begin{ques}\label{q1}
Mesons and Baryons and Pentaquarks (modified Griffiths 4.38, 4.35 in second edition)

Quarks carry total spin \( s = \frac{1}{2} \). Three quarks bind together to make a baryon (such as a proton or neutron), while two quarks (technically a quark and an antiquark, but the distinction is irrelevant here) bind together to make a meson (such as a pion or kaon). Assume the quark states are in the ground state with zero orbital angular momentum.
\begin{itemize}
\item[(a)] What spins are possible for baryons?

\item[(b)] What spins are possible for mesons?

\item[(c)] On July 3, 2015, the LHCb collaboration at CERN announced the discovery of a pentaquark state, consisting of \( 5 \) quarks. What spins are possible for a pentaquark?

\item[(d)] You can produce beams of baryons, mesons, and pentaquarks in your lab, and shoot them at Stern-Gerlach apparati. How many beams might you expect each of the samples to split into as it passes the SGA?
\end{itemize}
\end{ques}

\begin{proof}
    
    Here, we'll assume that under the tensor product, a spin-$\frac{n}{2}$ particle always has degeneracy $n+1$ associated (for $m=\{\frac{n}{2},\frac{n}{2}-1,...,-\frac{n}{2}+1, -\frac{n}{2}\}$). In a more mathematical language, as an irreducible $\so(3)$-representation, spin-$\frac{n}{2}$ corresponds to an $n$-dimensional irreducible representation.

    Also, we'll assume that each representation (for instance, the tensor representation, or sum of spins) is completely reducible, meaning it can be written as direct sums of various spin-$\frac{n}{2}$ representations.

    \hfil

    \begin{itemize}
        \item[(a)] For baryons that's bind by three quarks, the states are represented by tensor product of three spin-$\frac{1}{2}$ Hilbert spaces (each has dimension $2$), hence has dimension $2^3 = 8$. If consider the $S_u$ operator for $u\in\{x,y,z\}$, one has the following:
        \begin{align}
            S_u = S_{u,1}\tensor I_2\tensor I_3+I_1\tensor S_{u,2}\tensor I_3 + I_1\tensor I_2\tensor S_{u,3}
        \end{align}
        Which, the square is given as follow:
        \begin{align}
            S_u^2 =& S_{u,1}^2\tensor I_2\tensor I_3+I_1\tensor S_{u,2}^2\tensor I_3+I_1\tensor I_2\tensor S_{u,3}^2\\ 
            &+ 2(S_{u,1}\tensor S_{u,2}\tensor I_3+S_{u,1}\tensor I_2\tensor S_{u,3}+I_1\tensor S_{u,2}\tensor S_{u,3})
        \end{align}
        As a result, one has the operator $S^2$ given as follow:
        \begin{align}
            S^2 =& S_x^2+S_y^2+S_z^2\\
            =& S_1^2 \tensor I_2\tensor I_3+I_1 \tensor S_2^2\tensor I_3+I_1\tensor I_2\tensor S_3^2\\
            &+ 2\sum_{u=x,y,z}(S_{u,1}\tensor S_{u,2}\tensor I_3+S_{u,1}\tensor I_2\tensor S_{u,3}+I_1\tensor S_{u,2}\tensor S_{u,3})
        \end{align}
        Using the fact that $2\sum_{u=x,y,z}S_{u,1}\tensor S_{u,2} = S_{+,1}\tensor S_{-,2}+S_{-,1}\tensor S_{+,2}+2S_{z,1}\tensor S_{z,2}$ (derived for two spin-$\frac{1}{2}$ particle system), together with the associativity of tensor product (up to isomorphism), one has the following:
        \begin{align}
            S^2 =& S_1^2 \tensor I_2\tensor I_3+I_1 \tensor S_2^2\tensor I_3+I_1\tensor I_2\tensor S_3^2\\
            &+ S_{+,1}\tensor S_{-,2}\tensor I_3+S_{-,1}\tensor S_{+,2}\tensor I_3+2S_{z,1}\tensor S_{z,2}\tensor I_3\\
            &+ S_{+,1}\tensor I_2\tensor S_{-,3}+S_{-,1}\tensor I_2\tensor S_{+,3}+2S_{z,1}\tensor I_2\tensor S_{z,3}\\
            &+ I_1\tensor S_{+,2}\tensor S_{-,3}+I_1 \tensor S_{-,2}\tensor S_{+,3}+2I_1\tensor S_{z,2}\tensor S_{z,3}
        \end{align}
        Which, its action on the state $\ket{\uparrow}\tensor\ket{\uparrow}\tensor \ket{\uparrow}$ is as follow (Notice that if there is an $S_{+,i}$ involved in the tensor, then its action is $0$, since $S_{+,i}\ket{\uparrow_i} = \bzero$; on the other hand, each $S_i^2$ has action of scalar multiply by $\frac{3\hbar^2}{4}$, eigenvalue for spin-$\frac{1}{2}$, and each $S_{z,i}$ has action $\frac{\hbar}{2}$):
        \begin{align}
            S^2 \ket{\uparrow}\tensor\ket{\uparrow}\tensor \ket{\uparrow} = \left(3 \cdot \frac{3\hbar^2}{4} + 3 \cdot 2 \cdot \frac{\hbar^2}{4}\right)\ket{\uparrow}\tensor\ket{\uparrow}\tensor \ket{\uparrow} = \frac{15\hbar^2}{4}\ket{\uparrow}\tensor\ket{\uparrow}\tensor \ket{\uparrow}
        \end{align}
        Which, if $s$ satisfies $\hbar^2 s(s+1) = \frac{15\hbar^2}{4}$, then $s^2+s-\frac{15}{4} = 0$, or $(s+\frac{5}{2})(s-\frac{3}{2})=0$. Hence, for $s>0$, one has $s=\frac{3}{2}$. Which, it corresponds to degeneracy $4$, hence there is a spin-$\frac{3}{2}$ representation (with dimension $5$) in the tensor presentation, so there is an extra $8-4=4$-dimensional subrepresentation we're seeking in the tensor representation.

        \hfil

        For the extra $4$-dimensional subrepresentation, it's either spin-$\frac{3}{2}$ (with degeneray $4$), spin-$1$ and spin-$0$ direct sum (with degeneracy $3+1$), or spin-$\frac{1}{2}$ and spin-$\frac{1}{2}$ direct sum (with degeneracy $2+2$).

        Now,consider the operator again, if viewed the first two states as a composite state (i.e. do the tensor of the first two spaces first), using $S_{12}^2 = S_1^2\tensor I_2+I_1\tensor S_2^2 + S_{+,1}\tensor S{-,2}+S_{-,1}\tensor S_{+,2}+2 S_{z,1}\tensor S_{z,2}$ represents the spin operator, $S_{z,12}=S_{z,1}\tensor I_2+I_1 \tensor S_{z,2}$ represents its $z$-spin operator, and $S_{\pm,12}= S_{\pm,1}\tensor I_2+I_1 \tensor S_{\pm,2}$ represents the raising and lowering operator (which coincides with the definition $S_{x,12}\pm iS_{y,12}$ if work out the derivation), notice that it can be rewritten as follow:
        \begin{align}
            S^2 =& (S_1^2\tensor I_2+I_1\tensor S_2^2 + S_{+,1}\tensor S{-,2}+S_{-,1}\tensor S_{+,2}+2 S_{z,1}\tensor S_{z,2})\tensor I_3 + I_{12}\tensor S_3^2\\
            &+ (S_{+,1}\tensor I_2+I_1\tensor S_{+,2})\tensor S_{-,3}+(S_{-,1}\tensor I_2+I_1\tensor S_{-,2})\tensor S_{+,3}\\
            &+2(S_{z,1}\tensor I_2+I_1\tensor S_{z,2})\tensor S_{z,3}\\
            =& S_{12}^2\tensor I_3+I_{12}\tensor S_3^2 + S_{+,12}\tensor S_{-,3}+S_{-,12}\tensor S_{+,3} + 2S_{z,12}\tensor S_{z,3}
        \end{align}
        So, it can also be viewed as a composite (two spin-$\frac{1}{2}$) state, tensor a spin-$\frac{1}{2}$ state. If consider the state $\ket{\psi}:=\frac{\sqrt{2}\ket{1,1}\tensor \ket{-\frac{1}{2}}-\ket{1,0}\tensor \ket{\frac{1}{2}}}{\sqrt{3}}$ (for $\ket{s,m}$ being the spin-1 state in the composit of the first two, and $\ket{\pm \frac{1}{2}}$ represents spin up/down for the third particle), then $S^2$ has the following action:
        \begin{align}
            S^2\ket{\psi} =& \sqrt{\frac{2}{3}}\left[\left(2\hbar^2+\frac{3\hbar^2}{2}-\hbar^2\right)\ket{1,1}\tensor\ket{-\frac{1}{2}} + \sqrt{2}\hbar^2\ket{1,0}\tensor \ket{\frac{1}{2}}\right]\\
            &- \frac{1}{\sqrt{3}}\left[\left(2\hbar^2+\frac{3\hbar^2}{2}+0\right)\ket{1,0}\tensor\ket{\frac{1}{2}} + \sqrt{2}\hbar^2\ket{1,1}\tensor\ket{-\frac{1}{2}}\right]\\
            =& \frac{1}{\sqrt{3}}\left[\sqrt{2}\cdot\frac{3\hbar^2}{2}\ket{1,1}\tensor\ket{-\frac{1}{2}}-\frac{3\hbar^2}{2}\ket{1,0}\tensor\ket{\frac{1}{2}}\right]\\
            =& \frac{3\hbar^2}{2}\frac{\sqrt{2}\ket{1,1}\tensor\ket{-\frac{1}{2}}-\ket{1,0}\tensor\ket{\frac{1}{2}}}{\sqrt{3}}= \frac{3\hbar^2}{2}\ket{\psi}
        \end{align}
        This shows that $\ket{\psi}$ has spin $s$ satisfies $\hbar^2 s(s+1)=\frac{3\hbar^2}{2}$, which corresponds to spin $s=\frac{1}{2}$.
        This shows that there is a copy of spin-$\frac{1}{2}$ representation in the tensor product space of three spin-$\frac{1}{2}$ particles, which corresponds to degeneracy $2$.

        \hfil

        Eventually, one has spin-$\frac{3}{2}$ (degeneracy 4) and spin-$\frac{1}{2}$ (degeneracy 2) in an $8$-dimensional space, which the remaining $8-4-2=2$-dimension must be coming from spin-$\frac{1}{2}$ (the only spin with 2-dimensional irreducible $\so(3)$-representation). So, the possible spins for baryons are $s = \frac{3}{2},\frac{1}{2}$ (where there are $2$ copies of spin-$\frac{1}{2}$ when chosen a specific direct sum decomposition).

        \hfil

        \hfil

        \item[(b)] For mesons that's bind by two quarkes, the states are represented by tensor product of two spin-$\frac{1}{2}$ Hilbert space. Which, we've already calculated the sum of two spin-$\frac{1}{2}$, which has allowed spin $s=1,0$.
        
        \hfil

        \hfil

        \item[(c)] For a pentaquark state with 5 quarks bind together, it can be viewed as a system binding two smaller systems, one with 3 quarks binded (like a baryon, with Hilbert Space dimension $8$), the other with 2 quarks binded (like a meson, with Hilbert Space dimension $4$). So, for the two Hilbert Spaces, one has the first one with allowed spins $s_1 = \frac{3}{2},\frac{1}{2}$, and the second one with $s_2 = 1,0$.
        
        Let $\mathcal{H}^{(1)}$ denotes the first Hilbert Space, $\mathcal{H}^{(2)}$ denotes the second Hilbert Space, and $\mathcal{H}_s$ denotes the Hilbert Space with spin $s$. Based on the answer in part (a) and (b), one has the following:
        \begin{align}
            \mathcal{H}^{(1)} \cong \mathcal{H}_{\frac{1}{2}}\tensor\mathcal{H}_{\frac{1}{2}}\tensor \mathcal{H}_{\frac{1}{2}} \cong \mathcal{H}_{\frac{3}{2}}\oplus\mathcal{H}_{\frac{1}{2}}\oplus \mathcal{H}_{\frac{1}{2}},\quad \mathcal{H}^{(2)} \cong \mathcal{H}_{1}\oplus \mathcal{H}_{0}\cong \mathcal{H}_{\frac{1}{2}}\tensor \mathcal{H}_{\frac{1}{2}}
        \end{align}
        (Note: Reacll that composite two spin-$\frac{1}{2}$ particles have Hilbert Space $\mathcal{H}_{\frac{1}{2}}\tensor \mathcal{H}_{\frac{1}{2}}$; on the other hand, composite three spin-$\frac{1}{2}$ particles have Hilbert space $\mathcal{H}_{\frac{1}{2}}\tensor\mathcal{H}_{\frac{1}{2}}\tensor \mathcal{H}_{\frac{1}{2}}$).

        So, the pentaquark has its Hilbert Space (as an $\so(3)$-representation) being as follow:
        \begin{align}
            \mathcal{H}^{(1)}\tensor \mathcal{H}^{(2)} &\cong (\mathcal{H}_{\frac{3}{2}}\tensor \mathcal{H}_{1})\oplus (\mathcal{H}_{\frac{3}{2}}\tensor \mathcal{H}_{0})\oplus (\mathcal{H}_{\frac{1}{2}}\tensor\mathcal{H}_{\frac{1}{2}}\tensor \mathcal{H}_{\frac{1}{2}})\oplus (\mathcal{H}_{\frac{1}{2}}\tensor\mathcal{H}_{\frac{1}{2}}\tensor \mathcal{H}_{\frac{1}{2}})\\
            &\cong (\mathcal{H}_{\frac{3}{2}}\tensor \mathcal{H}_{1})\oplus (\mathcal{H}_{\frac{3}{2}}\tensor \mathcal{H}_{0})\oplus \mathcal{H}^{(1)}\oplus \mathcal{H}^{(1)}
        \end{align}
        First, let's resolve the case of $\mathcal{H}_0 \tensor \mathcal{H}_s$ for arbitrary $s$: Its spin operator $S^2 := S_1^2\tensor I_2+I_1\tensor S_2^2 + S_{+,1}\tensor S_{-,2}+S_{-,1}\tensor S_{+,2}+2S_{z,1}\tensor S_{z,2}$; but notie that $\mathcal{H}_0$ has everything acting on it trivially (since $s=m=0$), then anything involving nontrivial spin operators in the first column is a zero operator. Hence, $S^2 = I_1 \tensor S_2^2$, showing that $\mathcal{H}_0\tensor \mathcal{H}_s \cong \mathcal{H}_s$ (since $\mathcal{H}_0$ provides no action, and as a 1-dimension space it adds no dimension to the tensor product). Hence, the above tensor product can be reduced the following:
        \begin{align}
            \mathcal{H}^{(1)}\tensor \mathcal{H}^{(2)} \cong (\mathcal{H}_{\frac{3}{2}}\tensor \mathcal{H}_{1})\oplus \mathcal{H}_{\frac{3}{2}}\oplus \mathcal{H}^{(1)}\oplus \mathcal{H}^{(1)}
        \end{align}
        Now, since $\mathcal{H}^{(1)}$ has dimension $8$ (with spin $s=\frac{3}{2},\frac{1}{2}$), $\mathcal{H}_{\frac{3}{2}}$ has dimension $4$ (with spin $\frac{3}{2}$), with the tensor representation having dimension $8\times 4 = 32$, there are total of $32-8-8-4=12$-dimension remaining. Which all comes from $\mathcal{H}_{3/2}\tensor \mathcal{H}_1$.

        Doing all the manual calculation will put too much work for this, so we'll appeal to the Clebsch-Gordan table. Which, a composite spin of $\frac{3}{2}$ and $1$ particles will form allowed spin $s = \frac{5}{2},\frac{3}{2},\frac{1}{2}$ (with degeneracy $6,4,2$ respectively, matching up with the dimension $12$).

        Hence, here is the full direct sum decomposition:
        \begin{align}
            \mathcal{H}^{(1)}\tensor \mathcal{H}^{(2)} &\cong (\mathcal{H}_\frac{5}{2}\oplus \mathcal{H}_\frac{3}{2}\oplus \mathcal{H}_\frac{1}{2}) \oplus (\mathcal{H}_\frac{3}{2}) \oplus (\mathcal{H}_{\frac{3}{2}}\oplus\mathcal{H}_{\frac{1}{2}}\oplus \mathcal{H}_{\frac{1}{2}})\oplus (\mathcal{H}_{\frac{3}{2}}\oplus\mathcal{H}_{\frac{1}{2}}\oplus \mathcal{H}_{\frac{1}{2}})\\
            &= \mathcal{H}_\frac{5}{2}\oplus \mathcal{H}_\frac{3}{2}^{\oplus 4}\oplus \mathcal{H}_\frac{1}{2}^{\oplus 5}
        \end{align}
        Which, the allowed spin for the pentaquark is $s=\frac{5}{2},\frac{3}{2},\frac{1}{2}$.

        \hfil

        \item[(d)] If a beam contains baryons (spin $s=\frac{3}{2},\frac{1}{2}$ allowed), mesons (spin $s=1,0$ allowed), and pentaquarks (spin $s=\frac{5}{2},\frac{3}{2},\frac{1}{2}$ allowed), then the possible value of $m$ is $\{\frac{5}{2},\frac{3}{2},\frac{1}{2},-\frac{1}{2},-\frac{3}{2},-\frac{5}{2},1,0,-1\}$. So, there's possibly $9$ beams when passing through the Stern-Gerlach apparati.
    \end{itemize}
\end{proof}

\newpage

\section*{2}
\begin{ques}\label{q2}
Coupled Spins

Consider two particles, with \( s_1 = \frac{1}{2} \) and \( s_2 = 1 \). The particles are at rest and subject to an interaction that depends on the relative spin orientations, with Hamiltonian

\[
H = \epsilon \, \mathbf{S}_1 \cdot \mathbf{S}_2
\]

What are the energy eigenvalues for this Hamiltonian, and what are their degeneracies?
\end{ques}

\begin{proof}
    
\end{proof}

\newpage

\section*{3}
\begin{ques}\label{q3}
Clebsch Gordan Practice (this is like Griffiths 4.40, or 4.36 in second edition, if you want to practice on that first)
\begin{itemize}
\item[(a)] A particle of spin \( s_1 = \frac{3}{2} \) and a particle of spin \( s_2 = \frac{1}{2} \) are at rest in a configuration where the total spin of the two particles is \( 1 \) and a measurement of the total spin of the two particles in the \( z \) direction yields \( -\hbar \). If you measured the spin of the spin-\( \frac{3}{2} \) particle in the \( z \) direction, what values might you obtain, and with what probabilities?

\item[(b)] Imagine the outcome of your measurement in part (a) was \( -\frac{3\hbar}{2} \). You subsequently measure the total spin of the two-particle system. What values might you obtain and with what probabilities?

\item[(c)] An electron with spin up in the \( z \) direction is in the state \( \psi_{4,2,-1} \) of the hydrogen atom. If you measure the total angular momentum (orbital plus spin) what values might you obtain and with what probabilities?

\item[(d)] Imagine the outcome of your measurement of the total angular momentum in part (c) gave \( j = \frac{3}{2} \). You subsequently measure the spin of the electron in the \( z \) direction. What is the probability that your result is \( -\frac{\hbar}{2} \)?
\end{itemize}
\end{ques}

\begin{proof}
    
\end{proof}

\newpage

\section*{4 D}
\begin{ques}\label{q4}
EPR Pair (Griffiths 12.1 in both editions)

\hfil

The singlet spin configuration (State $\frac{1}{\sqrt{2}}(\ket{\uparrow\downarrow}-\ket{\downarrow\uparrow})$) is a classic example of an \emph{entangled state}--a two particle state that cannot be expressed as the product of two one-partile states, and for which, therefore, one cannot really speak of ``the state'' of either particle separately. You might wonder whether this is somehow an artifact of bad notation--maybe some linear combination of the one-particle states would disentangle the system. Prove the following theorem:
\begin{thm}
    Consider a two-level system, $\ket{\phi_a}$ and $\ket{\phi_b}$, with $\braket{\phi_i}{\phi_j}=\delta_{ij}$ (For example, $\ket{\phi_a}$ might represent spin up and $\ket{\phi_b}$ spin down). The two-particle state 
    $$\alpha\ket{\phi_a(1)}\ket{\phi_b(2)}+\beta\ket{\phi_b(1)}\ket{\phi_a(2)}$$
    with $\alpha,\beta\neq 0$ \emph{cannot} be expressed as a product $\ket{\psi_r(1)}\ket{\psi_s(2)}$ for \emph{any} one-particle states $\ket{\psi_r}$ and $\ket{\psi_s}$.
\end{thm}
\end{ques}

\begin{proof}
    Suppose the contrary that it can be expressed as a product $\ket{\psi_r(1)}\tensor\ket{\psi_s(2)}$ for some other one-particle states $\ket{\psi_r}$ and $\ket{\psi_s}$. If express them as linear combinations of $\ket{\phi_a}$ and $\ket{\phi_b}$, one gets the following for some $x,y,z,w\in\CC$ satisfying $|x|^2+|y|^2=|z|^2+|w|^2=1$:
    \begin{align}
        \ket{\psi_r} =x\ket{\phi_a}+y\ket{\phi_b},\quad \ket{\psi_s}=z\ket{\phi_a}+w\ket{\phi_b} 
    \end{align}
    Then, its tensor product is expressed as follow:
    \begin{align}
        \ket{\psi_r}\tensor\ket{\psi_s} = xz\ket{\phi_a(1)}\tensor\ket{\phi_a(2)}+xw\ket{\phi_a(1)}\tensor\ket{\phi_b(2)}+yz\ket{\phi_b(1)}\ket{\phi_a(2)}+yw\ket{\phi_b(1)}\tensor\ket{\phi_b(2)}
    \end{align}
    Which, compared to the original state $\ket{\psi_r}\ket{\psi_s}=\alpha\ket{\phi_a(1)}\ket{\phi_b(2)}+\beta\ket{\phi_b(1)}\ket{\phi_a(2)}$, since under the tensor product, the pure tensor of the basis elements form a basis for the tensor product, each coefficient must match up. Hence, we get the following:
    \begin{align}
        xz = yw = 0,\quad xw = \alpha,\quad yz = \beta
    \end{align}
    This indicates that either $x$ or $z$ is $0$, and also either $y$ or $w$ is $0$. Notice that in each case, at least one of $\alpha,\beta$ is $0$:
    \begin{itemize}
        \item If $x=0$, then $\alpha=xw=0$
        \item If $z=0$, then $\beta=yz=0$
        \item If $y=0$, then $\beta=yz=0$ again
        \item If $w=0$, then $\alpha=xw=0$ again.
    \end{itemize}
    This contradicts the assumption that $\alpha,\beta\neq 0$.

    So, it can't be written in a pure tensor form of any vectors in the space (i.e. it's not a product of any one-particle states).
\end{proof}

\newpage

\section*{5 D}
\begin{ques}\label{q5}
\( S_x \) of the singlet state

The singlet state

\[
|\psi_S\rangle = \frac{1}{\sqrt{2}} \left( |\uparrow\downarrow\rangle - |\downarrow\uparrow\rangle \right)
\]

has \( s = 0 \), and therefore any measurement of \( S^2 \), \( S_x \), \( S_y \), or \( S_z \) will give zero. In this problem, show explicitly that

\[
S^2 |\psi_S\rangle = 0
\]

and that

\[
S_x |\psi_S\rangle = 0 .
\]

The most straightforward way to approach the second part of this problem is to re-write \( |\psi_S\rangle \) in terms of the \( S_x \) eigenstates that were derived in the text book or in class.
\end{ques}

\begin{proof}
    First, recall that for the addition of two spin-$\frac{1}{2}$ particles has the following definition for $S^2$ and $S_x$:
    \begin{align}
        &S^2 = S_1^2 \tensor I_2 + I_1 \tensor S_2^2 + (S_{1,+}\tensor S_{2,-}+S_{1,-}\tensor S_{2,+}) + 2(S_{1,z}\tensor S_{2,z})\\
        &S_x = S_{x,1}\tensor I_2 + I_1\tensor S_{x,2}
    \end{align}    
    Hence, we get the following:
    \begin{align}
        S^2\ket{\psi_S} =& \frac{1}{\sqrt{2}}S^2(\ket{\uparrow}\tensor\ket{\downarrow}-\ket{\downarrow}\tensor\ket{\uparrow})\\
        =& \frac{1}{\sqrt{2}}\left(\bigg(\frac{3\hbar^2}{4} + \frac{3\hbar^2}{4}-\frac{\hbar^2}{2}\bigg)\ket{\uparrow}\tensor\ket{\downarrow}+\hbar^2\ket{\downarrow}\tensor\ket{\uparrow}\right) - \frac{1}{\sqrt{2}}\left(\left(\frac{3\hbar^2}{4}+\frac{3\hbar^2}{4}-\frac{\hbar^2}{2}\right)\ket{\downarrow}\tensor\ket{\uparrow}+\hbar^2\ket{\uparrow}\tensor\ket{\downarrow}\right)\\
        =& \frac{\hbar^2}{\sqrt{2}}(\ket{\uparrow\downarrow}+\ket{\downarrow\uparrow})- \frac{\hbar^2}{\sqrt{2}}(\ket{\uparrow\downarrow}+\ket{\downarrow\uparrow})\\
        =& \bzero
    \end{align} 
    \begin{align}
        S_x\ket{\psi_S} &= \frac{1}{\sqrt{2}}(S_{x,1}\tensor I_2+I_1\tensor S_{x,2})(\ket{\uparrow}\tensor\ket{\downarrow}-\ket{\downarrow}\tensor\ket{\uparrow})\\
        &= \frac{1}{\sqrt{2}}\left(\left(\frac{\ket{\uparrow}+\ket{\downarrow}}{\sqrt{2}}\right)\tensor\ket{\downarrow} - \left(\frac{\ket{\uparrow}-\ket{\downarrow}}{\sqrt{2}}\right)\tensor\ket{\uparrow} + \ket{\uparrow}\tensor\left(\frac{\ket{\uparrow}-\ket{\downarrow}}{\sqrt{2}}\right)-\ket{\downarrow}\tensor\left(\frac{\ket{\uparrow}+\ket{\downarrow}}{\sqrt{2}}\right)\right)\\
        &= \frac{1}{2}(\ket{\uparrow}\tensor\ket{\downarrow}+\ket{\downarrow}\tensor\ket{\downarrow}-\ket{\uparrow}\tensor\ket{\uparrow}+\ket{\downarrow}\tensor\ket{\uparrow}+\ket{\uparrow}\tensor\ket{\uparrow}-\ket{\uparrow}\tensor\ket{\downarrow}-\ket{\downarrow}\tensor\ket{\uparrow}-\ket{\downarrow}\tensor\ket{\downarrow})\\
        &= \bzero
    \end{align}
\end{proof}

\newpage

\section*{6 D}
\begin{ques}\label{q6}
Is it entangled?

Is

\[
|\psi\rangle = \frac{1}{2} \left( |\uparrow\uparrow\rangle + |\uparrow\downarrow\rangle - |\downarrow\uparrow\rangle + |\downarrow\downarrow\rangle \right)
\]

entangled? Prove it one way or the other.
\end{ques}

\begin{proof}
    Given the above states, we'll prove that it's entangled. 

    \hfil
    
    Suppose the contrary that it's not entangled, it can be written as $\ket{\theta_1}\tensor\ket{\phi_2}$ a pure tensor of two states in the 2-level system. Rewrite each $\ket{\theta},\ket{\phi}$ under the basis $\ket{\uparrow},\ket{\downarrow}$, there exists $x,y,z,w in \CC$ with $|x|^2+|y|^2=|z|^2+|w|^2=1$, such that the following is true:
    \begin{align}
        \ket{\theta} = x\ket{\uparrow}+y\ket{\downarrow},\quad \ket{\phi}=z\ket{\uparrow}+w\ket{\downarrow}
    \end{align}
    Then, the pure tensor becomes the following:
    \begin{align}
        \ket{\psi} = \ket{\theta_1}\tensor\ket{\phi_2} = xz\ket{\uparrow}\tensor\ket{\uparrow}+xw\ket{\uparrow}\tensor\ket{\downarrow}+yz\ket{\downarrow}\tensor\ket{\uparrow}+yw\ket{\downarrow}\tensor\ket{\downarrow}
    \end{align}
    To matchh up with the original expression $\ket{\psi}=\frac{1}{2}(\ket{\uparrow}\tensor\ket{\uparrow}+\ket{\uparrow}\tensor\ket{\downarrow}-\ket{\downarrow}\tensor\ket{\uparrow}+\ket{\downarrow}\tensor\ket{\downarrow})$, one must have the following:
    \begin{align}
        xz = xw = yw = \frac{1}{2},\quad yz = -\frac{1}{2}
    \end{align}
    As a result, one has $x,y,z,w\neq 0$. So, the eqution $xz = xw$ implies $z=w$; yet, this also implies that $yz = yw = -\frac{1}{2}$, which contradicts the equation $yw = \frac{1}{2}$.

    So, this state can't be written as a pure tensor of any states, showing $\ket{\psi}$ has the two particles being entangled.
\end{proof}

\end{document}

